\documentclass{article}

% ready for submission
%\usepackage[accepted]{icml2026}

\usepackage[utf8]{inputenc} % allow utf-8 input
\usepackage[T1]{fontenc}    % use 8-bit T1 fonts
\usepackage{hyperref}       % hyperlinks
\usepackage{url}            % simple URL typesetting
\usepackage{booktabs}       % professional-quality tables
\usepackage{amsfonts}       % blackboard math symbols
\usepackage{nicefrac}       % compact symbols for 1/2, etc.
\usepackage{microtype}      % microtypography
\usepackage{xcolor}         % colors
\usepackage{svg}
\newcommand{\theHalgorithm}{\arabic{algorithm}}

\usepackage{icml2026}


% From ICML
\usepackage{amsmath}
\usepackage{amssymb}
\usepackage{mathtools}
\usepackage{amsthm}
\usepackage{dsfont}
\usepackage[most]{tcolorbox}
\usepackage{microtype}
\usepackage{graphicx}
\usepackage{subfigure}

\usepackage{csquotes}

\usepackage[capitalize,noabbrev]{cleveref}

\usepackage{tabularx}
\newcolumntype{Y}{>{\centering\arraybackslash}X}
\usepackage{multirow}

%%%%%%%%%%%%%%%%%%%%%%%%%%%%%%%%
% THEOREMS
%%%%%%%%%%%%%%%%%%%%%%%%%%%%%%%%
\theoremstyle{plain}
\newtheorem{theorem}{Theorem}[section]
\newtheorem{proposition}[theorem]{Proposition}
\newtheorem{lemma}[theorem]{Lemma}
\newtheorem{corollary}[theorem]{Corollary}
\theoremstyle{definition}
\newtheorem{definition}[theorem]{Definition}
\newtheorem{assumption}[theorem]{Assumption}
\theoremstyle{remark}
\newtheorem{remark}[theorem]{Remark}


% custom from us
\usepackage{enumitem}
\usepackage{xcolor}
\usepackage{colortbl}
\definecolor{lightgray}{gray}{0.85}
\definecolor{darkgray}{gray}{0.7}
%\definecolor{customblue}{rgb}{0.8, 0.9, 1}
\definecolor{customblue}{rgb}{0.27, 0.36, 0.55}

% Todonotes is useful during development; simply uncomment the next line
%    and comment out the line below the next line to turn off comments
\usepackage[disable,textsize=tiny]{todonotes}
\usepackage{array}
\usepackage{makecell}
%\usepackage[textsize=tiny]{todonotes}

\input{math_commands.tex}

\newlength{\myboxwidth}
\setlength{\myboxwidth}{7cm}
\newcommand{\myfont}{
    %\large       % Font size
    \ttfamily    % Typewriter (monospace) font
}


\begin{document}

\twocolumn[

\icmltitle{\texorpdfstring{Position: We Need to Realign Incentives for Meaningful Progress\\in Adversarial Alignment for LLMs}{Realigned Incentives for Adversarial Alignment}}

% It is OKAY to include author information, even for blind
% submissions: the style file will automatically remove it for you
% unless you've provided the [accepted] option to the icml2025
% package.

% List of affiliations: The first argument should be a (short)
% identifier you will use later to specify author affiliations
% Academic affiliations should list Department, University, City, Region, Country
% Industry affiliations should list Company, City, Region, Country

% You can specify symbols, otherwise they are numbered in order.
% Ideally, you should not use this facility. Affiliations will be numbered
% in order of appearance and this is the preferred way.
\icmlsetsymbol{equal}{*}

\begin{icmlauthorlist}
\icmlauthor{Anonym}{}


%\icmlauthor{}{sch}
%\icmlauthor{}{sch}
\end{icmlauthorlist}

\icmlaffiliation{Anonym}{Anonym}
\icmlcorrespondingauthor{Anonym}{Anonym}

% You may provide any keywords that you
% find helpful for describing your paper; these are used to populate
% the "keywords" metadata in the PDF but will not be shown in the document
\icmlkeywords{Machine Learning, ICML, LLM robustness, judging, autograder}

\vskip 0.3in
]

% this must go after the closing bracket ] following \twocolumn[ ...

% This command actually creates the footnote in the first column
% listing the affiliations and the copyright notice.
% The command takes one argument, which is text to display at the start of the footnote.
% The \icmlEqualContribution command is standard text for equal contribution.
% Remove it (just {}) if you do not need this facility.

%\printAffiliationsAndNotice{}  % leave blank if no need to mention equal contribution
\printAffiliationsAndNotice{\icmlEqualContribution} % otherwise use the standard text.

\begin{abstract}
Misaligned research objectives have considerably hindered progress in adversarial robustness research over the past decade.
For instance, an extensive focus on optimizing target metrics, while neglecting rigorous standardized evaluation, has led researchers to pursue ad-hoc heuristic defenses that were seemingly effective. Yet, most of these were exposed as flawed by subsequent evaluations, ultimately contributing little measurable progress to the field.
In this position paper, we illustrate that current research on the robustness of large language models (LLMs) risks repeating past patterns with potentially worsened real-world implications. To address this, we argue that \textbf{realigned objectives are necessary for meaningful progress in adversarial alignment.}
To this end, we build on established cybersecurity taxonomy to formally define differences between past and emerging threat models that apply to LLMs.
Using this framework, we illustrate that progress requires disentangling adversarial alignment into addressable sub-problems and returning to core academic principles, such as measurability, reproducibility, and comparability.
Although the field presents significant challenges, the fresh start on adversarial robustness offers the unique opportunity to build on past experience while avoiding previous mistakes.
\end{abstract}


\section{Introduction}

%\tom{also what seems quite different is that the input space bound of the so-called toy problem (Lp norm balls) were mostly for semantic indifference. now we dont need any constraint on the input since a semantic change can be classified in the correct answer. so the input constraint can be moved to the output classification.}

Security risks in computer science have been a prevalent issue for decades~\cite{valiant_learning_1985, kearns_learning_1993}.
This ongoing challenge has resulted in an ``arms race" that includes the continued development of new attacks, such as malware and phishing, as well as defense mechanisms. %Here, an arms race refers to the ongoing competition among multiple entities, each trying to improve upon the other by employing progressively advanced techniques to either breach or safeguard a specific computer system.

In the field of deep learning, \citet{szegedy_intriguing_2014} discovered that deep neural networks are highly susceptible to \textit{adversarial examples} -- input perturbations optimized to mislead models into making predictions that are erroneous or misaligned with their intended behavior. In response, countless defense strategies were proposed to safeguard neural networks against these attacks in the last decade. Yet, most newly proposed heuristic defenses were exposed as flawed by subsequent evaluations, often by using standard attack protocols that already existed at the time of the defense publication~\citep{tramer_adaptive_2020}. 

We investigate the research state of adversarial alignment in large language models (LLMs), which we define as the ability of an aligned model to maintain its intended training objective in the presence of adversarial attacks. 
We observe that adversarial alignment risks following the same cycle of flawed defenses and subsequent rectified evaluations seen in past adversarial robustness research but with considerably amplified challenges and stakes~\cite{hendrycks2022x, zou2023universal, schwinn2023adversarial}. 
%For instance, reliable evaluation is hindered by increased problem complexity, such as larger model scale as well as semantic and discrete input and output domains~\cite{andriushchenko2024jailbreaking, li2024llm}. 
Unlike previous robustness research that generally focused on well-defined problems like image classification~\cite{goodfellow_explaining_2015}, assessing LLM capabilities is considerably more challenging due to inherent ambiguities of the alignment problem and the complexity of natural language~\cite{wolf2023fundamental, andriushchenko2024jailbreaking, li2024llm}. Furthermore, the potential harm associated with LLMs is substantially greater due to their advanced capabilities and widespread availability~\cite{hendrycks2022x} in particular as they start being used as autonomous agents.  %This makes slow research progress and resulting misaligned and non-robust models a serious concern.

%\leo{Sophie/Cas/Leo: ML research is benchmark driven. We even advocate for it (leader boards clearly defined metrics than can be improved). But must fulfil certain criteria (e.g., the ones we advocate reproducibility, comparability, measureability). In the adversarial ML setting it also should be as easy to reliable evaluate as possible (e.g., no defense through obfuscation). Important to distangle the problem (ill defined measures of progress) and the thing we propose (well defined problem definitions, simpler problems etc.)}

\begin{table*}[t!]
    \centering
    \caption{Non-exhaustive comparison of robustness research in previous domains and LLMs.}
     %\resizebox{1\textwidth}{!}{
     \newcommand{\cellwidth}{2.2cm}
     \newcommand{\linesp}{3pt}
     \newcommand{\cbox}[1]{\parbox[t][1cm][t]{\cellwidth}{\centering #1}}
    %\setlength{\arrayrulewidth}{10pt}
    \tiny	
\begin{tabular}{l*{3}{>{\centering\arraybackslash}m{\cellwidth}}*{3}{>{\centering\arraybackslash}m{\cellwidth}}}
    \addlinespace
     & \multicolumn{3}{@{\hspace{3mm}}c@{\hspace{3mm}}}{\cellcolor{customblue}\textcolor{white}{\textbf{Attack goals}}} & \multicolumn{3}{@{\hspace{3mm}}c@{\hspace{3mm}}}{\cellcolor{customblue}\textcolor{white}{\textbf{Attack Capabilities}}} \\
    \addlinespace[\linesp]
    \rowcolor{darkgray} & \textbf{Objectives} & \textbf{Attacks} & \textbf{Datasets} & \textbf{Access} & \textbf{Constraints} & \textbf{Frameworks} \\
    \rowcolor{lightgray}\textbf{Previous} & \cbox{Clear objectives\\(e.g., classification)\\\cite{szegedy_intriguing_2014, madry_towards_2018}} & \cbox{Generally reliable\\\cite{carlini_evaluating_2019, tramer_adaptive_2020, croce2020reliable}} & \cbox{Standardized\\\cite{croce2020reliable, Schwinn2021Jitter, croce2020robustbench}} & \cbox{Generally white-box\\and open-source\\\cite{szegedy_intriguing_2014, croce2020robustbench}} & \cbox{Tractable but incomplete\\(e.g., $\ell_p$)~\cite{szegedy_intriguing_2014, goodfellow_explaining_2015, madry_towards_2018}} & \cbox{Standardized\\\cite{croce2020robustbench, papernot2018cleverhans, rauber2017foolbox}} \\
    \addlinespace[\linesp]
    \rowcolor{lightgray}\textbf{LLM} & \cbox{Alignment \& robustness entangled~\cite{zou2023universal, mazeika2024harmbench}} & \cbox{Currently weak\\\cite{andriushchenko2024jailbreaking, li2024llm}} & \cbox{Entangled notions of harmfulness\\\cite{mazeika2024harmbench, chao2024jailbreakbench}} & \cbox{Often black-box and\\proprietary models\\\cite{zou2023universal, chao2023jailbreaking}} & \cbox{No constraints\\\cite{zou2023universal, mazeika2024harmbench, chao2023jailbreaking, zhu2023autodan}} & \cbox{Varying evaluation settings\\\cite{mazeika2024harmbench, chao2024jailbreakbench}} \\
    %\bottomrule
\end{tabular}
    %}
    \label{tab:comparison}
    \vspace{-15pt}
\end{table*}

We argue that the past lack of progress can be largely attributed to a narrow focus on improving benchmark numbers without sufficient attention to rigorous evaluations and clear evaluation criteria. This led to the proliferation of ad-hoc defenses that relied on security through obscurity and ultimately proved ineffective, thus not providing a solid foundation for future work. Now, the field of adversarial alignment risks repeating the same mistakes, expending significant effort but failing to make meaningful progress.  
%it repeats past mistakes with larger-scale consequences. 
%\,\,\textbf{1)} vague problem definitions, where alignment and robustness are inherently entangled in the adversarial alignment problem, \textbf{2)} complex and non-reproducible evaluations, and \textbf{3)}  an emphasis on state-of-the-art attack performance against proprietary models over reproducible and comparable open-source research.
Our position on adversarial alignment in LLMs is as follows:

\begin{tcolorbox}[
    colback=white, colframe=customblue, coltitle=white, fonttitle=\bfseries, 
    rounded corners, enhanced, 
    title=Position, 
    attach boxed title to top left={yshift=-2mm, xshift=5mm}, 
    boxed title style={colback=customblue, rounded corners},
    boxsep=1mm,
    left=1mm, % adjust left padding
    right=1mm
]
 %Current research incentives in LLMs:\,\,\textbf{1)} encourage vague problem definitions and complex, non-reproducible evaluations; \textbf{2)} promote a focus on SOTA performance on proprietary models over reproducible and comparable open-source research.
 We argue that researchers deal with:\,\,\textbf{1)} vague problem definitions, where alignment and robustness are inherently entangled in the adversarial alignment problem, \textbf{2)} complex, non-reproducible evaluations, and \textbf{3)} emphasis on state-of-the-art attack performance against proprietary models over reproducible and comparable open-source research. \textbf{Thus, meaningful progress in adversarial alignment for LLMs requires simpler, reproducible, and more measurable objectives.}
\end{tcolorbox}

Our main contributions are:


% entangled definitions of robustness in current benchmarks --> Simpler datasets with narrowly defined simple threat settings
% entangled definition of adeversarial alignment (alignment and robustness) require complex evaluation pipelines (e.g., judges) and can lead to ambigituies  --> Proxy objectives
% evals on proprietary models are incentivized to show real-world consequence of attack --> hard to reproduce / compare / measure --> open-source evals should be preferred in most settings and for other cases research focused APIs could improve reproducibility
% Community leaderboards and standatarized evals

%\renewcommand{\labelitemi}{--}
\begin{itemize}[noitemsep,nolistsep,topsep=-2pt,leftmargin=0.6cm] 
    \item We systematically identify challenges in previous robustness research, how they apply to LLMs, and how new challenges emerged. Based on this analysis: 
%\end{itemize}
%Based on this analysis:
%\begin{itemize}[noitemsep,nolistsep,topsep=-2pt,leftmargin=0.6cm] 
    \item We demonstrate how adversarial alignment intertwines the challenges of alignment and robustness, making robustness evaluation difficult, as it inherits the challenges of measuring alignment, such as ambiguity in success criteria. 
    We advocate for a complementary approach that directs the majority of technical works toward simpler sub-problems with measurable objectives, while reserving different criteria for exploratory research.
    
    \item Towards the same goal of improving measurability, we propose simplifying robustness benchmarks by evaluating specific types of harm individually rather than combining complex concerns like copyright infringement, fairness, and toxicity into a single evaluation.
    
    \item We emphasize the need for academia to prioritize reproducible research over chasing SOTA performance on proprietary models. In this context, we propose fostering open-source research with accessible models. 
    
    \item Computational overhead, vast number of hyperparameters, and varying implementation details hinder comparability between different works. We advocate for a practical approach to improve reproducibility and comparability: community-driven leaderboards and standardized benchmarks to encourage best practices in adversarial robustness research.  
\end{itemize}

For a concise discussion and supporting literature regarding the claims made in this work, we refer the reader to Appendix~\ref{app:evidence}.

%\item We advocate for evaluating attack effectiveness using predefined optimization objectives rather than relying on subjective LLM-judged outputs, increasing methodological transparency.  

%Our goal is to highlight that current research risks repeating past patterns. We illustrate that an impending arms race entails considerable challenges that, if not addressed by the community, could result in misaligned incentives that hinder research efforts, slow down the adversarial alignment of LLMs, and have severe real-world consequences. 

%We adopt a cybersecurity taxonomy, encompassing robustness goals, and attack capabilities, to analyze the evolving threat landscape for large language models (LLMs). This analysis reveals how the increasing complexity of LLMs, coupled with these misaligned incentives, risks repeating past failures. We specifically highlight how a disproportionate focus on narrow attack vectors and easily achievable defense benchmarks has hindered the development of robust LLMs. To address these challenges and enable meaningful progress in adversarial alignment, we advocate for a fundamental shift in research incentives. We posit that academia, by concentrating on foundational research, pursuing unique ideas, and generating novel insights while upholding the core principles of reproducibility, comparability, and rigorous measurability, can play a pivotal role in this transformation. This requires a concerted effort to reduce problem complexity, such as disentangling robustness from other concerns, to create a research environment that fosters genuine advancements in LLM robustness.

% \begin{figure*}
%     \centering
%     \includegraphics[width=1\linewidth]{figures/f1-cropped.pdf}
%     \caption{Preliminary - illustrate paper structure}
%     \label{fig:structure}
% \end{figure*}

\section{Structure of our argument}\label{sec:structure}

We structure our position as follows: We define a taxonomy of adversarial robustness threats based on common cybersecurity frameworks including \textbf{I)} robustness goals and \textbf{II)} adversary capabilities. For both elements of this taxonomy, we follow a parallel argument structure: We \textbf{A)} define past and present threat models (see start of \S\ref{sec:pos-goals} and \S\ref{sec:pos-capability}), \textbf{B)} discuss historical challenges and connections to upcoming and exacerbated issues in the LLM domain (\S\ref{sec:pos-goals-challenges}, and \S\ref{sec:pos-capability-challenges}), and \textbf{C)} explore the applicability of past insights to new problems and how current research objectives can be realigned to promote measurable progress (\S\ref{sec:pos-goals-realigned} and \S\ref{sec:pos-capability-realigned}). We provide a non-exhaustive comparison between past and current robustness research in Table~\ref{tab:comparison}. (Appendix~\ref{app:evidence}). \\
Additionally, Appendix~\ref{app:other_works} presents a comprehensive comparison between our work and other contributions from the AI safety community, including differences to the concurrently submitted position paper by~\citet{rando2025adversarial}. %Overall, we structure our position to emphasize the novel research challenges emerging in adversarial alignment for LLMs and provide a forward-looking perspective.

% in this field. 

%We propose that new objectives should facilitate core academic values including reproducibility (amount of effort needed to reproduce an experiment), comparability (even reproducible experiments may be difficult to compare across works), and measurability (metrics should be quantitative and robust to design choices). We argue that such objectives reinforce academia’s strengths: exploring unique foundational ideas while maintaining methodological rigor. Unlike industry, which excels at scaling, academia, constrained by resources, is uniquely positioned to generate reliable insights through open, transparent, and reproducible research. Moreover, its independent incentive structure -- akin to white-hat hackers in cybersecurity -- enables rigorous testing of AI systems without being tied to product-driven and comercial incentives.
%2) These incentives should reward academia’s unique role in generating unconventional ideas. By enabling independent researchers to explore diverse directions, we build an ``idea pipeline'' for industry labs to evaluate and scale. This creates an environment where academia’s creativity meets industry’s capacity for large-scale development, driving scientific and technological progress forward.

% \section{Related Positions and Future Outlook} 

% \textbf{Concurrent positions.} Our findings are reinforced by the concurrent work of~\citet{rando2025adversarial}, who independently identify similar fundamental challenges in adversarial alignment of LLMs. 
% Their research offers detailed case studies of specific threat models like jailbreaks, poisoning, and unlearning. While we investigate specific failures as well--albeit in less detail--, our focus is on developing a high-level taxonomy and proposing forward-looking solutions to emerging challenges. 
% %Their paper offers detailed case studies of specific threat models like jailbreaks, poisoning, and unlearning, whereas our work develops a high-level taxonomy and focuses more on forward-looking solutions. 
% Together, both works offer complementary perspectives toward a more comprehensive understanding of the identified issues.

% \textbf{Call for positions.} We will revisit and evaluate the arguments presented in this position paper one year or more from its publication. This retrospective analysis will examine whether our identified challenges and proposed solutions have proven relevant regarding the field's progress in adversarial robustness for LLMs. We warmly welcome other researchers to contribute their perspectives and collaborate on the retrospective, whether they want to support or challenge the positions outlined in the remainder of this paper. Through this collaborative reflection, we aim to maintain accountability for our arguments.

%The concurrent position paper by \citet{rando2025adversarial} emphasizes the importance of properly tackling adversarial alignment. While our work provides a high-level taxonomy and focuses more on forward-looking solutions, \citet{rando2025adversarial} perform a more detailed investigation of specific threat models like jailbreaks, poisoning, and unlearning.


%\textbf{Concurrent Position.} Concurrent to our work,~\citet{rando2025adversarial} take a similar position and argue that adversarial alignment in LLMs studies problems that are vaguely defined, more difficult to evaluate, and likely harder to solve. They also emphasize the need for caution to ensure these challenges do not impede overall research progress in the field.

%Case studies \\
%High-level forward looking
\input{parts/taxonomy}
\input{parts/main_paper.tex}

%\section*{Impact Statement}\label{app:impact}
%Adversarial attacks on LLMs can have considerable consequences in real-world applications. Still, as machine-learning robustness has been an unsolved research problem for the last decade, we believe that the best way to approach this problem is through culminating awareness. Currently, it seems unlikely that the robustness issue can be completely resolved through technical means. Thus, making people aware of the harmful use cases and limitations of these models appears to be necessary to avoid irresponsible deployment of such models for critical applications and to reduce the harm malicious actors can cause. Moreover, this work does not introduce new technical innovations. Instead, it addresses impediments in robustness research and explores potential solutions to accelerate research progress.


\bibliographystyle{icml2026}
\bibliography{references}

%%%%%%%%%%%%%%%%%%%%%%%%%%%%%%%%%%%%%%%%%%%%%%%%%%%%%%%%%%%%

\clearpage

\appendix

\section{Appendix}


\subsection{Related work}\label{app:other_works}

Here, we discuss the differences to other works on AI safety and one related position paper that was concurrently submitted to our work.

\subsubsection{Difference to other works related to AI safety}

The development of safe and reliable AI, particularly Large Language Models (LLMs), is a central concern. As a result, there are numerous works which address the importance of reliable and comprehensive safety evaluations in the LLM safety community~\cite{bommasani2021opportunities, weidinger2021ethical, ganguli2022red, liang2022holistic, hendrycks2023overview, srivastava2022beyond, chen2024trustworthy, bai2022constitutional, shevlane2023model, bowman2022measuring, sillberg2024eu, casper2023open, anwar2024foundational, casper2024black, schwinn2023adversarial, casper2023explore, reuel2024open}.

Yet, all of these contributions have substantially different goals compared to the proposed position. One group of works primarily evaluate capabilities and outline potential risks~\cite{bommasani2021opportunities, weidinger2021ethical, hendrycks2023overview, srivastava2022beyond, liang2022holistic, shevlane2023model, anwar2024foundational}, another is focused on methodology and technical limitations and advancements~\cite{ganguli2022red, chen2024trustworthy, bai2022constitutional, bowman2022measuring, casper2023open, casper2023explore}, while others focus on formulating policy recommendations required for safe AI usage now and in the future~\cite{sillberg2024eu, casper2024black, reuel2024open}. Our work, in contrast, focuses specifically on the process of research itself within adversarial alignment, diagnosing how current practices risk hindering cumulative progress by not sufficiently learning from the history of adversarial robustness research.

In the following we provide a categorization of prior work from the AI safety community. Note, that some work could be assigned to multiple categories.

\textbf{Discussions on AI capabilities and risks.} A significant body of work explores the capabilities of LLMs and the spectrum of risks they introduce. For instance, \citet{bommasani2021opportunities} provide a broad overview of foundation models capabilities and resulting risks, while \citet{weidinger2021ethical} detail potential ethical and social harms stemming from language models. The potential for catastrophic risks is further analyzed in \citet{hendrycks2023overview, shevlane2023model}. Other research highlights areas of concern giving potential future capabilities of LLMs and discusses current limitations \cite{srivastava2022beyond}. More recently, foundational challenges in ensuring LLM alignment and safety have been systematically documented, pointing to the deep-seated difficulties in achieving these goals \cite{anwar2024foundational}.

\textbf{Methodological improvements for AI safety.} The LLM safety community has proposed numerous methodological advancements. This includes novel alignment techniques, such as Constitutional AI, which aims to instill harmlessness through AI feedback \cite{bai2022constitutional}. Red teaming has emerged as a critical methodology for discovering vulnerabilities \cite{ganguli2022red, casper2023explore}, and significant effort is dedicated to measuring progress on scalable oversight \cite{bowman2022measuring}. Methodological critiques also exist, for example, highlighting the limitations of RLHF \cite{casper2023open}. 

\textbf{AI policy and governance recommendations.} Recognizing the societal impact of AI, researchers and institutions are increasingly proposing policy and governance frameworks. These range from critiques and recommendations concerning specific regulations like the EU AI Act \cite{sillberg2024eu}, to broader architectural frameworks for trustworthy and responsible AI \cite{chen2024trustworthy}, and discussions on open problems in technical AI governance \cite{reuel2024open}. Such works aim to guide the development and deployment of AI in a manner that aligns with societal values and minimizes harm. Moreover, they advocate for safety policy, where technical solutions alone might fail. 

In summary, existing literature offers critical insights into AI risks, develops novel methodologies, and proposes governance structures. Our work complements these efforts by adopting a historical and methodological lens focused on how to \textbf{facilitate research progress}. By drawing explicit lessons from the trajectory of adversarial robustness research in other domains, we identify current pitfalls and advocate for a realignment of research objectives to ensure that the field makes efficient and durable progress, avoiding the cycles of ineffective solutions that characterized earlier efforts in adversarial machine learning~\cite{madry_towards_2018, athalye_obfuscated_2018, tramer_adaptive_2020, schwinn2023adversarial}. This specific focus distinguishes our contribution.

\subsubsection{Concurrent position}

Our findings are reinforced by the concurrent work of~\citet{rando2025adversarial}, who independently identify similar fundamental challenges in adversarial alignment of LLMs. 
Their research offers detailed case studies of specific threat models like jailbreaks, poisoning, and unlearning. We also investigate specific failures as--albeit in less detail. Moreover, in contrast to their work, we develope a high-level taxonomy encompassing all sorts of threat models. Furthermore, we go beyond problem identification to propose actionable, forward-looking solutions that address emerging challenges in adversarial alignment research.


% %Their paper offers detailed case studies of specific threat models like jailbreaks, poisoning, and unlearning, whereas our work develops a high-level taxonomy and focuses more on forward-looking solutions. 
% Together, both works offer complementary perspectives toward a more comprehensive understanding of the identified issues.


\subsection{Further arguments against the position}~\label{app:arguments_against}

Another argument against a focus on well-defined sub-problems is the broad attack surface of foundation models from chat assistants to autonomous agents. The majority of threats related to unexplored and upcoming applications will not be formally well-defined. Here, research requires identifying novel threat surfaces in new application areas. This kind of research has been proven to be highly relevant in the past, including the identification of vulnerabilities in commercial hashing algorithms~\cite{levenson2021apple}, seemingly reliable protections of art against generative AI~\cite{honig2024adversarial}, or vulnerabilities of newly developed benchmarks, such as LLM leaderboard~\cite{huang2025exploring}.

\textbf{New opportunities for defenses.} In this work, we primarily emphasize the increased complexity of achieving robustness in LLMs compared to previous research domains. However, the semantic output and input space of LLMs not only complicate measuring robustness but also open up new opportunities for defenses. Early adversarial detection methods in computer vision focused on identifying \textit{small}, \textit{imperceptible} norm-bounded perturbations, proved largely ineffective~\cite{carlini2017adversarial}. Moreover, detection in the output space was not possible due to the prevalence of simple classification tasks. Conversely, recent advancements in the LLM domain demonstrate the strong empirical results of input and output classifiers for defending against harmful requests~\cite{kim2023robust, sharma2025constitutional}. This suggests that the inherent semantics of the LLMs domain, while presenting new challenges, also facilitate the development of new defense mechanisms previously unavailable in other domains.

\textbf{Capability progress is too fast.} Another argument against the proposed position of reliable research in constrained settings is the rapid pace of progress in model capabilities. By simplifying the problem, researchers risk creating solutions that will quickly become obsolete as AI models evolve. This could be specifically the case as current adversarial attack algorithms often rely on LLMs to jailbreak other models. Rather than trying to solve isolated, small problems in the field, researchers should concentrate on large, fundamental questions that are not dependent on current technical limitations, as they may be resolved with progress in LLM capability. Moreover, while open-source models are increasing in their capabilities, research on open-source models may always lag behind proprietary ones and academia may recover mistakes already identified by industry labs. 

\subsection{Beyond Discussed Threats}~\label{app:beyond}
Most arguments presented in this work regarding emerging challenges in LLM robustness research are applicable beyond the specific threat models discussed as part of the taxonomy. 
These include issues related to complex input-output domains (e.g., precisely defining attack constraints or evaluation goals) and challenges associated with large model sizes, vast number of hyperparameters, and varied implementation settings.
This includes integrity threats, such as model poisoning, as well as attacks on confidentiality (e.g., membership inference, model inversion) and hybrid threats, like model unlearning. 
However, we note that these related areas are likely to present their own unique complexities and challenges within the LLM context that are currently hindering research progress, which requires discussion in future work.

\subsection{Examples of Measurability Improvements in Recent Work}\label{app:measurability}
Here, we highlight recent datasets that demonstrate how complex problems can be made more measurable and list desirable properties they fulfill:
\begin{itemize}
    \item \textbf{Jailbreak Tax}~\cite{nikolic2025jailbreak}: Multiple choice outcomes, verifiable outcomes.
    \item \textbf{AutoAdvExBench}~\cite{carlini2025autoadvexbench}: Quantifiable metrics, real-world coding task.
    \item \textbf{NYU CTF Bench}~\cite{shao2025nyu}: Unit-testable, cyber security challenges.
    \item \textbf{AgentHarm}~\cite{andriushchenko2024agentharm}: Verifiable actions.
\end{itemize}

Beyond highlighting desirable properties in existing benchmarks, we propose three ideas to improve measurability:
\begin{enumerate}
    \item \textbf{Robust judgeability}: For open-ended datasets such as JailbreakBench or HarmBench, curate subsets of ``robustly judgeable examples'' with low judge disagreement across defenses, models, and attacks.
    \item \textbf{Direct dataset construction}: Directly construct datasets with straightforward success criteria, such as highly toxic completions identifiable by keywords.
    \item \textbf{Guaranteed success settings}: Define settings where an attack is guaranteed to succeed in principle. Following the approach proposed in~\cite{carlini2023aligned}, we identify extremely unlikely output sequences for a model via brute-force search and optimize the attack to elicit them. This ensures that the adversarial optimization problem can be solved and studied in isolation.
\end{enumerate}

\subsection{Evidence for Claims}\label{app:evidence}
In this section, we provide a more detailed breakdown of the central claims presented in this position paper. For each claim, we synthesize evidence from the literature to support our position and propose concrete experimental directions to further validate our claims in the future.

\begin{enumerate}
    \item \textbf{Claim:} \textit{Complex, open-ended evaluations often result in measurement errors, whereas measurable sub-problems provide more reliable signals of progress.}
    
    \textbf{Evidence:} 
    Evaluating the success of an adversarial attack on a broad, open-ended task (e.g., "generate a harmful response") relies heavily on subjective interpretation and complex automated judges, which introduces significant noise~\cite{beyer2025llm}. Recent benchmarks demonstrate that narrowing the domain improves reliability. For instance, the \textbf{Jailbreak Tax}~\cite{nikolic2025jailbreak} reveals that open-ended evaluations often systematically overestimate attack success due to high false-positive rates in LLM judges. By restricting the evaluation to a "measurable sub-problem", such as a multiple-choice format or verifiable domains (e.g., math), researchers can isolate the performance of the attack algorithm from the vagaries of the evaluation pipeline. Similarly, \textbf{NYU CTF Bench}~\cite{shao2025nyu} and \textbf{AgentHarm}~\cite{andriushchenko2024agentharm} show that unit-testable challenges (e.g., capturing a flag or performing a specific function call) offer unambiguous success criteria, eliminating the need for subjective judgment and improving measurability.
    
    \textbf{Proposed Experiment:} 
    Future studies could correlate performance improvements on these measurable, narrow tasks with performance on broader, open-ended datasets. If improvements in optimization on narrow tasks consistently transfer to open-ended settings, it would validate the use of measurable sub-problems as reliable proxies for general progress.

    \item \textbf{Claim:} \textit{Open-ended evaluations suffer from significant variability in judge performance, undermining comparability.}
    
    \textbf{Evidence:} 
    The reliance on "LLM-as-a-Judge" for evaluating open-ended adversarial attacks introduces a critical source of instability. Studies such as those by \citet{li2025llms} and \citet{beyer2025llm} highlight substantial inconsistencies between different LLM judges, as well as discrepancies between LLM judges and human annotators. These judges can be sensitive to irrelevant formatting details or biases in the prompt, leading to false positives that vary across different attack methods. This variability means that a \enquote{successful} defense in one paper might fail under a slightly different judge configuration in another, making it impossible to reliably compare defense mechanisms.
    
    \textbf{Proposed Experiment:} 
    Extend existing studies to systematically investigate differences in false positive/negative rates of various judges specifically across different defense and attack algorithms. Quantifying how much the choice of judge alters the \textit{ranking} of defenses would demonstrate the severity of this issue.

    \item \textbf{Claim:} \textit{Reliance on proprietary model evaluations creates significant reproducibility risks.}
    
    \textbf{Evidence:} 
    Evaluations conducted on proprietary APIs (e.g., GPT, Claude) are inherently unstable due to unannounced updates, version deprecations, and opaque system changes. A defense or attack evaluated on "GPT-5.2" today may yield completely different results next week if the backend model or safety filter is updated. This \enquote{drifting baseline} makes it impossible to reproduce results from prior literature or to benchmark progress over time accurately. The deprecation of specific API snapshots further exacerbates this, rendering older experimental setups permanently inaccessible.
    
    \textbf{Proposed Experiment:} 
    Conduct a longitudinal study tracking the success rate of a fixed set of attacks against proprietary model APIs over a period of 6-12 months. Quantifying the variance in results solely due to API changes would provide concrete evidence of the \enquote{reproducibility crisis} in proprietary evaluations.

    \item \textbf{Claim:} \textit{Research progress relies on the adoption of common, standardized benchmarks.}
    
    \textbf{Evidence:} 
    The history of adversarial robustness in computer vision demonstrates the necessity of standardization. The creation of \textbf{RobustBench}~\cite{croce2020robustbench} provided a unified framework that allowed the community to rigorously track progress, filter out ineffective defenses, and identify genuine state-of-the-art methods. In contrast, the current state of LLM robustness research is fragmented, with many papers using custom datasets, evaluation protocols, and threat models~\cite{beyer2025llm}. This fragmentation obscures whether new methods are genuinely better or simply optimized for a specific, non-standard setting.
    
    \textbf{Proposed Experiment:} 
    Encourage and track the adoption of third-party evaluations for LLM defenses, similar to the community effort seen in RobustBench. A study analyzing the "breakage rate" of defenses when evaluated on a standardized benchmark versus their reported performance would highlight the gap caused by lack of standardization.

    \item \textbf{Claim:} \textit{The performance gap between open-source and proprietary models has narrowed, making open-source models viable proxies for research.}
    
    \textbf{Evidence:} 
    Historically, proprietary models were significantly more capable than open-source alternatives, justifying their use in evaluations. However, recent trends show a rapid convergence. Leaderboards like the \textbf{LMSYS Chatbot Arena} show open-weights models (e.g., GLM-4.7, Kimi-k2-Thinking-Turbo, and DeepSeek-V3.2) performing on par with or even exceeding frontier proprietary systems. This suggests that insights gained from robustifying or attacking these strong open models are likely to transfer to proprietary counterparts, without the reproducibility downsides of API-based research. 
    
    \textbf{Proposed Experiment:} 
    Track whether methodological trends (e.g., susceptibility to specific attack vectors, effectiveness of defenses) observed on state-of-the-art open models are predictive of trends in proprietary models. If the correlation is high, it validates the shift towards open-source-first research.

    \item \textbf{Claim:} \textit{Academic research on open models is a driver of safety improvements in proprietary systems.}
    
    \textbf{Evidence:} 
    Despite the focus on proprietary models in some benchmarks, foundational methodological advances in AI safety often originate in academia using open models. Key examples include the \textbf{GCG} attack~\cite{zou2023universal}, the \textbf{HarmBench} framework~\cite{mazeika2024harmbench}, and early adversarial training algorithms~\cite{xhonneux2024efficient}. These methods, developed on open weights, have heavily influenced the safety reports and defense strategies of major proprietary labs (e.g., referenced in Gemini and LLaMA technical reports), demonstrating that open-model research directly impacts the wider ecosystem.
    
    \textbf{Proposed Experiment:} 
    Perform a citation and genealogy analysis of safety techniques deployed in major industrial models to quantify the proportion of core safety methodologies that originated from academic research on open systems.

    \item \textbf{Claim:} \textit{Current automated attacks overestimate robustness compared to human adaptive attacks.}
    
    \textbf{Evidence:} 
    Automated attack algorithms, while improving, still struggle with the complex optimization landscape of language. Evidence from recent competitions and studies~\cite{li2024llm, andriushchenko2024jailbreaking, nasr2025attacker} shows that manual \enquote{jailbreaking} or human-in-the-loop strategies consistently achieve higher Attack Success Rates (ASR) on robust models than fully automated baselines. This gap implies that relying solely on current automated attacks for evaluation can lead to a false sense of security, as the model may be vulnerable to a more creative or adaptive human adversary.
    
    \textbf{Proposed Experiment:} 
    Benchmark state-of-the-art automated attacks against a coordinated human red-teaming effort on the same set of defended models at the end of 2026. Quantifying the \enquote{automation gap} would emphasize the need for stronger adaptive attack algorithms (or proxies) before automated evaluation can be fully trusted.
\end{enumerate}


\end{document}